\chapter{Background and Related Work}
\label{ch:background}

Chapter~\ref{ch:introduction} identified three requirements for sensing-assisted
vehicular beam alignment: distributed multimodal perception that can reveal
relevant physical targets, target-to-user association that assigns sensing
evidence to the correct communication entity, and a resource-aware evaluation
that determines whether the resulting beam-management policy is worthwhile.
This chapter reviews the literature through those requirements rather than as a
collection of disconnected sensing modalities or learning architectures. The
review first establishes the measurement burden in high-frequency V2I beam
management and the opportunity offered by sensing-aided beam prediction. It
then considers distributed multimodal BEV perception, target-to-user association
with temporal tracking, and system-level V2I evaluation. The chapter concludes
by positioning the three contributions of this thesis relative to the respective
boundaries of prior work.

\section{Beam Management and Sensing-Aided Beam Prediction}
\label{sec:background-beam-management}

\subsection{Measurement Burden in High-Frequency V2I Communication}
\label{subsec:background-codebooks-training}

High-frequency V2I communication relies on directional transmission to obtain
sufficient array gain. A BS consequently has to identify a beam that is suitable
for each active vehicle and revise that decision as vehicle geometry, visibility,
and propagation conditions evolve. In practical codebook-based systems, this
process is supported by initial-access sweeps, reference-signal measurements,
and later refinement. The measurements are valuable because they are tied
directly to the communication link, but they consume radio resources and become
less informative as the vehicle moves or the environment changes.

The resulting trade-off is especially acute on roads. Relative position changes
quickly, vehicles and roadside structures intermittently block dominant paths,
and several users may occupy similar angular sectors. A policy that measures a
large beam set frequently can maintain reliable link evidence, but it reduces
the resources available for data transmission. A policy that measures a small
set risks missing a useful beam or retaining stale channel information. This
relationship between beam-search cost and communication performance motivates
beam prediction as a candidate-reduction mechanism rather than as an isolated
classification task \cite{xu2022computervisionaidedmmwave,liu2020radarassistedpredictivebeamformingvehicular,liu2022learningbasedpredictivebeamformingintegrated}.

Evaluation of beam-management methods should therefore distinguish several
quantities. A beam-ranking metric indicates whether a short candidate list
contains a useful direction, whereas a resource metric records the measurements
required to obtain that list. Neither quantity alone describes the communication
consequence. The latter also depends on how the selected beam affects the
effective channel, how often additional recovery measurements are required, and
how the measurement resources are accounted for. This distinction is important
when interpreting the sensing-aided studies reviewed below.

\subsection{From Sensing-Aided Prediction to Candidate Reduction}
\label{subsec:background-sensing-aided-prediction}

Sensing-aided beam prediction exploits the fact that wireless propagation is
coupled to the physical scene. Early vision-aided work showed that cameras at a
BS can assist blockage prediction and beam selection \cite{alrabeiah2019millimeterwavebasestations}. Subsequent V2I studies used visual information to infer beam
alignment or beam-coherence behaviour for moving vehicles
\cite{xu2022computervisionaidedmmwave}. Position-assisted methods use coarse
geometric information, while radar- and LiDAR-aided methods use range, angular,
and motion cues that are more directly linked to vehicle geometry
\cite{charan2021visionpositionmultimodalbeamprediction,demirhan2021radaraided6gbeam,jiang2022lidaraidedfuturebeam}. These approaches demonstrate that environment
observations can reduce the search space that must be examined by communication
measurements.

The prediction target varies across this literature. Some methods select a
single beam index, whereas others return a ranked candidate list or estimate the
power associated with each codebook beam. A distributional output is particularly
useful for beam management because it exposes both the ordering of candidates
and the concentration of the model's evidence. It can support Top-\(K\) candidate
selection, confidence assessment, and a later decision about whether the
candidate list is safe to use. Common metrics include Top-\(K\) accuracy,
received-power loss, and beam-search overhead. The DeepSense 6G data collection
and its associated beam-prediction challenge helped establish these metrics for
multimodal sensing-and-communication studies \cite{alkhateeb2023deepsense6glargescalerealworld,charan2022multimodalbeampredictionchallenge}.

Recent work extends this direction through multimodal transformers, sequential
models, and decision-oriented beam-management strategies
\cite{tian2023multimodaltransformerswirelesscommunications,ghassemi2024multimodaltransformerreinforcementlearningbased,farzanullah2025beamselectionisacusing}. These studies strengthen the case for using rich
environmental context, but they also reveal an important modelling boundary.
The learning sample is often already associated with a particular communication
link, either through the data construction or through an available position
label. In a multi-vehicle road scene, however, a BS first observes physical
objects and only subsequently has to determine which object is the UE represented
by a communication measurement. A useful beam ranking is therefore necessary,
but not sufficient, for user-specific beam alignment.

\section{Distributed Multimodal BEV Perception for Joint Vehicle Detection and Beam Prediction}
\label{sec:background-distributed-bev-perception}

\subsection{BEV Representation and Multimodal Fusion}
\label{subsec:background-bev-fusion}

Bird's-eye-view perception provides a natural spatial interface for roadside
multimodal sensing. Perspective camera images can be transformed into a
ground-plane representation, while point-cloud features can be encoded in the
same coordinate frame. Once expressed in BEV, information from different views
and modalities can be fused before target-level reasoning is performed. The
Lift-Splat-Shoot formulation illustrated how features from arbitrary camera rigs
can be lifted and projected into a BEV representation \cite{philion2020liftsplatshootencoding}. BEVFusion further showed the value of using a shared BEV
space to combine camera and geometric sensing for multiple perception tasks
\cite{liu2024bevfusionmultitaskmultisensorfusion}.

These ideas are useful for vehicular communication because beam quality is
strongly related to the geometry between the target BS, the vehicle, and the
surrounding environment. A BEV feature map can encode a vehicle's position,
local context, and relative visibility in the same representation from which
beam-related features are derived. Point-based detection methods such as
CenterNet offer an efficient way to represent vehicles by their centres and
associated attributes \cite{zhou2019objectspoints}. In a communication-oriented
extension, the target feature at a detected vehicle can also support prediction
of a beam-energy distribution. Thus, joint vehicle detection and beam prediction
need not be regarded as unrelated tasks: both depend on an accurate local scene
representation.

The transfer of BEV methods from autonomous driving to communication should,
however, be made carefully. Conventional BEV perception is commonly evaluated
through detection, segmentation, or motion forecasting. These outputs do not
automatically yield a communication decision. In the present context, the value
of BEV fusion lies in generating target-level evidence that can support a later
beam-management procedure. The perception output remains associated with a
physical target; it does not itself provide a communication UE identity or
confirm the final service beam.

\subsection{Distributed Sensing and Multimodal Beam Prediction}
\label{subsec:background-distributed-perception}

Cooperative perception extends single-site sensing by combining observations
from multiple agents or infrastructure nodes. In road environments, distinct
viewpoints can enlarge coverage, reduce the impact of an individual occlusion,
and make the scene representation less dependent on one camera's field of view.
V2X-ViT, for example, studies feature fusion across heterogeneous V2X agents
and highlights issues such as pose error and asynchronous observations
\cite{xu2022v2xvitvehicletoeverythingcooperativeperception}. These issues are equally relevant when a target BS draws on
neighbouring sensing nodes, even though the target BS remains the local
communication decision point.

Distributed sensing-aided networking and multimodal beam prediction provide the
closest communication-oriented precedents. Imran \emph{et al.} use distributed
RGB sensing nodes to provide environment semantics for beam prediction
\cite{imran2024environmentsemanticcommunicationenabling}. Multimodal prediction
studies combine image, point-cloud, radar, or position information to exploit
their complementary strengths \cite{tian2023multimodaltransformerswirelesscommunications,orimogunje2026occlusionawaremultimodalbeamprediction,mollah2026multimodalsensingfusionmmwave}. Zeng \emph{et al.} further demonstrate the potential of a BEV-fusion
formulation for sequential multimodal beam prediction
\cite{zeng2026bevfusionbasedframeworksequential}. Together, these studies show
that multimodal and spatially distributed observations can improve the evidence
available for beam-related inference.

Their scopes nevertheless leave an opportunity for a target-BS-centered
formulation that jointly treats roadside perception and user-level beam
management. Distributed sensing studies often exchange visual semantics rather
than fusing multiview cameras and BS-local ISAC-like geometry in a shared local
BEV frame. BEV beam-prediction studies do not necessarily distinguish a detected
physical target from the communication UE that will later consume the beam hint.
The first contribution of this thesis addresses this interface by using
distributed multimodal BEV perception for joint vehicle detection and per-target
beam prediction. Its purpose is to produce a structured physical-target
representation, not to bypass the association and measurement stages required
for beam alignment.

\section{Target-to-User Association with IMM-Based Tracking}
\label{sec:background-target-user-association-tracking}

\subsection{From Physical Targets to Communication Users}
\label{subsec:background-target-user-association}

The association problem arises because perception and communication attach
identity to different objects. A perception module detects physical vehicles,
and a tracking system assigns them temporal track identifiers. A communication
system, by contrast, schedules UEs and obtains link measurements indexed by
their communication identities. These identities may correspond to the same
vehicle, but the correspondence is not observed directly by a sensing model.
The distinction becomes material whenever multiple vehicles are visible, a
vehicle is temporarily missed, or an inactive object resembles an active UE.

Multi-candidate beam-prediction research begins to address this ambiguity.
Charan \emph{et al.} consider camera-based prediction in multi-candidate V2I
scenarios and introduce transmitter identification as part of the problem
\cite{charan2023camerabasedmmwavebeam}. Distributed environment-semantic
communication likewise includes transmitter identification and object
association-based tracking \cite{imran2024environmentsemanticcommunicationenabling}. These studies establish that object-level inference cannot be assumed to be
user-level inference. They do not, however, fully resolve how a target-to-user
association should be anchored by communication measurements, maintained after
temporary perception failures, and evaluated through its downstream influence on
beam-management decisions.

Target-to-user association in this thesis is therefore defined as a user-level
data-association problem. Sensing supplies a target-level beam likelihood and
communication measurement supplies UE-level beam evidence. A one-to-one
assignment connects the two only when the evidence is compatible and otherwise
preserves a conventional communication path. This interpretation differs from
using an oracle identifier, a geometric label, or a transmitter class as a
substitute for association. It establishes the correspondence needed for a beam
hint to be used by the correct UE.

\subsection{IMM-Based Tracking and Association Persistence}
\label{subsec:background-imm-tracking}

Association must also persist over time. Tracking-by-detection methods associate
new observations with previous tracks using spatial proximity, motion
consistency, appearance, or learned displacement. CenterTrack exemplifies a
point-based formulation that jointly reasons about detection and temporal
association \cite{zhou2020trackingobjectspoints}. In vehicular communication,
state prediction is additionally valuable because an anticipated vehicle motion
can inform when a previously established beam relationship should be refreshed.

ISAC predictive-beamforming studies provide complementary evidence for this
temporal viewpoint. Radar-assisted and Bayesian approaches use sensing-derived
state information to predict future beamforming decisions
\cite{liu2020radarassistedpredictivebeamformingvehicular,yuan2020bayesianpredictivebeamformingvehicular}. Learning-based methods use historical channels or sensing observations to capture
temporal dynamics for predictive beamforming \cite{liu2022learningbasedpredictivebeamformingintegrated,liu2022deepclstmpredictivebeamforming,du2021integratedsensingcommunicationsv2i}. The common premise is that vehicle state evolves continuously even when a
communication measurement is intermittent.

An interacting multiple-model (IMM) tracker is suitable for this setting because
it combines motion hypotheses that describe different vehicle behaviours, such
as approximately constant motion and manoeuvring turns. Its role in the present
thesis is deliberately specific. The IMM tracker maintains the temporal
validity of an association between a perception track and a communication UE;
it does not infer a UE identity from visual appearance. When the predicted track
and a new detection cease to agree, the resulting loss of confidence is a reason
to invoke measurement-assisted re-association rather than to continue applying
an old beam hint. The second contribution of this thesis consequently joins
target-to-user association with IMM-based tracking as a single mechanism for
identity-consistent beam alignment.

\section{Multimodal V2I Simulation and 5G NR-Oriented Beam-Management Evaluation}
\label{sec:background-v2i-evaluation}

\subsection{Multimodal Data and System-Level Simulation}
\label{subsec:background-multimodal-data}

Evaluation of sensing-assisted communication requires observations and channel
information that can be related at the same decision epoch. Multimodal data
collections such as DeepSense 6G provide synchronised communication and sensing
records that support controlled beam-prediction studies
\cite{alkhateeb2023deepsense6glargescalerealworld}. Realistic simulation
frameworks offer a complementary route by combining road traffic, environmental
sensing, and wireless propagation models. For example, Park \emph{et al.}
consider a multimodal realistic simulation framework for resource-efficient
beam prediction \cite{park2025resourceefficientbeampredictionmmwave}. These
resources make it possible to inspect how perception errors and communication
measurements interact, rather than assessing a learning model only on isolated
labels.

The Multimodal V2I Simulator considered in this thesis is a system-level
evaluation environment built on existing multimodal V2I data assets. It does not
claim to regenerate all sensing or propagation assets. Instead, it organises the
available observations, trajectories, beam-quality labels, and channel-related
information into beam-management episodes centred on the serving BS. This scope
is sufficient to compare the consequences of different sensing-assisted policies
while keeping the underlying data and radio assumptions common across methods.

\subsection{Resource-Aware Beam-Management Evaluation}
\label{subsec:background-resource-evaluation}

The evaluation of a sensing-assisted policy should preserve the distinction
between a prediction and a communication decision. A high Top-\(K\) beam metric
indicates that a small candidate set may contain a useful beam, but it does not
establish whether the beam was associated with the right UE, whether the
candidate set was actually measured, or whether the saved resources improved
the user experience. Conversely, an evaluation that reports an achievable rate
without accounting for the reference-signal resources can conceal the cost of
maintaining the required channel knowledge.

ISAC and learning-based beam-management studies have examined communication
objectives such as received power, spectral efficiency, sum rate, regret, and
latency \cite{liu2022learningbasedpredictivebeamformingintegrated,Wang_2023,ghassemi2024multimodaltransformerreinforcementlearningbased,farzanullah2025beamselectionisacusing,mollah2026multimodalsensingfusionmmwave}. Such objectives are necessary but must be interpreted under a stated
measurement model. A 5G NR-oriented abstraction provides this discipline by
representing the distinct roles of initial access, beam refinement, candidate
measurement, and data transmission without claiming to reproduce every
protocol-level procedure. It allows policies to be compared under a common
accounting of measurement resources and under common physical-layer assumptions.

The third contribution of this thesis combines the Multimodal V2I Simulator with
such a 5G NR-oriented beam-management evaluation framework. The evaluation
traces the error chain from target availability and beam ranking, through
target-to-user association and IMM-based tracking, to candidate recall,
measurement overhead, effective rate, lower-tail rate, and outage. It therefore
does not presuppose that an improvement in perception accuracy automatically
produces an improvement in communication performance. Rather, it makes the
conditions under which sensing evidence yields a favourable rate--overhead
trade-off explicit.

\section{Research Positioning and Summary}
\label{sec:background-gap}

The reviewed literature establishes the components required for sensing-assisted
beam alignment, but typically examines them at different points in the
sensing-to-communication chain. Sensing-aided beam-prediction studies show that
camera, radar, LiDAR, position, and multimodal observations can reduce the
uncertainty of beam selection. BEV fusion and cooperative perception show how
multi-view and multi-agent observations can form a more complete spatial
representation. Multi-candidate prediction, object association, and predictive
beamforming show that identity and temporal state matter once several moving
vehicles share the scene. Finally, system-level communication studies establish
the importance of rate and resource objectives. The remaining task is to connect
these elements without assuming away the target-to-user association problem.

The first contribution of this thesis is positioned at the interface between
distributed multimodal perception and communication-oriented inference. It
uses a target-BS-centered BEV representation to jointly detect vehicles and
predict their beam-energy distributions, thereby producing target-level evidence
that can reflect complementary viewpoints and sensing modalities. The second
contribution addresses the identity interface left open by target-level
prediction. It uses communication-side evidence for target-to-user association
and IMM-based tracking to maintain that association over time, so that a beam
hint is evaluated as user-specific information rather than as an anonymous
object prediction.

The third contribution completes the chain at the system level. The Multimodal
V2I Simulator and 5G NR-oriented beam-management evaluation framework compare
candidate-measurement and fallback decisions under common data and resource
assumptions. This framing makes it possible to determine whether a beam hint is
not only plausible for a visible target but also correctly attributed, measured,
and beneficial for communication. The thesis is thus positioned as an
integration of distributed multimodal BEV perception, target-to-user association
with IMM-based tracking, and resource-aware beam-management evaluation. The
next chapter formalises the corresponding network entities, observations,
identities, and decision variables.
